\begin{abstract}
Large-scale unsupervised language models (LMs) are adept at acquiring extensive world knowledge and basic reasoning abilities, but precisely directing their behavior is challenging because their training is entirely unsupervised. Common approaches to enhance such controllability involve gathering human judgments about the comparative quality of model outputs and fine-tuning the LM to conform to these preferences, frequently utilizing reinforcement learning from human feedback (RLHF). Yet, RLHF involves a cumbersome and sometimes unstable pipeline: first, constructing a reward model to emulate human preferences, then applying reinforcement learning to the original LM to increase the expected reward, all while constraining deviations from the base model. In this work, we present an alternative formulation for the reward model within RLHF that permits the analytic derivation of the optimal policy, enabling the RLHF task to be resolved with a straightforward classification objective. We term this method \textit{Direct Preference Optimization} (DPO), which offers a stable, efficient, and resource-friendly process, eliminating the necessity for sampling from the LM during fine-tuning or for extensive hyperparameter search. Our empirical results demonstrate that DPO fine-tunes LMs to agree with human preferences on par with, or superior to, previous techniques. In particular, DPO not only surpasses PPO-based RLHF in guiding model sentiment but also matches or outperforms it in summarization and single-turn dialogue tasks, all while being much easier to deploy and train.
\end{abstract}

\section{Introduction}
Large-scale unsupervised language models (LMs) trained on vast amounts of data have demonstrated remarkable emergent abilities~\citep{chowdhery2022palm, brown2020language, touvron2023llama,bubeck2023sparks}. Nevertheless, such models are developed using data created by humans who possess diverse motivations, priorities, and levels of expertise. Not all of these attributes are appropriate for AI models to replicate; for instance, although it is valuable for an AI-based programming assistant to \textit{recognize} common coding mistakes for the purpose of providing corrections, we would still prefer the assistant to emulate the rare but exemplary programming skills also present in its training set when generating new code. Likewise, although a language model ought to be \textit{informed} about widely held misconceptions—such as one believed by half the population—it should not represent this misunderstanding as fact half of the time it is queried about the topic. This highlights the importance of aligning the model’s \emph{target outputs and behaviors} with desired standards, as opposed to the full range of its \textit{acquired knowledge and capabilities}, in order to construct AI systems that are robust, reliable, and amenable to human direction~\citep{ouyang2022training}. Current approaches generally shape LMs to better accord with human judgments via reinforcement learning (RL); as we will discuss, however, the RL-based preference objective employed by these techniques can in fact be solved precisely through a straightforward binary cross-entropy loss, thus streamlining the preference optimization process.

Broadly, contemporary approaches instill target behaviors within language models by leveraging curated datasets that capture human evaluations of safe and helpful outputs. This phase of preference learning follows an initial stage of unsupervised pre-training on extensive text corpora. While supervised fine-tuning on human-annotated exemplary outputs represents the most direct strategy for preference alignment, the most widely adopted family of methods is reinforcement learning from human or artificial feedback (RLHF/RLAIF; \citep{christiano2017deep,bai2022constitutional}). RLHF begins by fitting a reward model to annotated human preference data, after which RL is employed to adjust the language model’s policy so that its responses attain higher predicted reward without significant divergence from the original policy. Despite the effectiveness of RLHF in yielding language models with strong dialogue and coding performance, this framework introduces a higher level of complexity compared to purely supervised methods, requiring the training of multiple LMs and in-loop policy sampling, which leads to increased computational demands.

This work demonstrates a method for optimizing language models to align with human preferences without resorting to explicit reward modeling or reinforcement learning. We introduce \textit{{Direct Preference Optimization} (DPO)}, an approach that indirectly maximizes the same objective function targeted by RLHF methods—namely, reward maximization under a KL-divergence constraint—while remaining both easy to implement and efficient to train. Conceptually, DPO modifies the log odds between favored and unfavored model outputs, applying a per-example importance weighting that guards against the type of model collapse observed when using a naive probability ratio objective. Like prior methods, DPO depends on a theoretical preference model such as the Bradley-Terry model \cite{bradley1952rankanalysis}, which assesses the degree to which a reward function agrees with observed preference judgments. Distinctively, whereas existing approaches leverage the preference model to train a reward function before optimizing the policy with respect to this learned reward, DPO leverages a variable transformation that expresses the preference loss directly as a function of the policy. As a result, given human preference data over model completions, DPO is able to train policies directly using a straightforward binary cross entropy loss, yielding an optimal policy corresponding to an implicit reward function fitted to these preferences.

Our primary contribution is Direct Preference Optimization (DPO): a simple algorithm, free of reinforcement learning, that facilitates language model training from preference feedback. Through empirical evaluation, we find that DPO matches or outperforms established techniques—including PPO-based RLHF—across tasks such as sentiment transformation, summarization, and conversational modeling, when applied to language models containing up to 6B parameters.

\section{Related Work}

Large self-supervised language models are capable of handling some tasks without explicit task-specific examples \citep{radford2019language} and can further improve with the use of few-shot prompting techniques \citep{gpt3,megatron,chowdhery2022palm}. Nevertheless, directly training these models on curated instruction and human-generated response datasets—an approach known as instruction-tuning—substantially boosts their effectiveness on various downstream applications and enhances their alignment with user intentions \citep{mishra-etal-2022-cross,sanh2022multitask,chung2022scaling,thoppilan2022lamda}. Through instruction-tuning, LLMs gain the capacity to handle instructions beyond those presented during fine-tuning, thereby increasing their generalizability and overall utility \citep{chung2022scaling}. 

In spite of the progress made through instruction tuning, it is often more feasible to obtain \textit{comparative} human evaluations of model output than it is to gather expert-crafted demonstrations. As a result, recent efforts have shifted toward refining LLMs using datasets reflecting human preferences, leading to improvements in tasks such as translation \citep{kreutzer-etal-2018-reliability}, summarization \citep{stiennon2022learning,ziegler2020finetuning}, narrative generation \citep{ziegler2020finetuning}, and instruction following \citep{ouyang2022training,ramamurthy2023is}. This methodology typically involves two steps: first, training a neural reward model to reflect human preferences using a preference-learning framework like the Bradley-Terry model \citep{bradley1952rankanalysis}; and second, optimizing the language model to maximize this reward function through reinforcement learning, often via REINFORCE \citep{williams1992reinforce}, proximal policy optimization (PPO; \cite{schulman2017proximal}), or related algorithms \citep{ramamurthy2023is}. 

Related research uses LLMs, already instruction-tuned with human feedback, to generate large-scale synthetic preference datasets that target qualities such as safety or harmlessness \citep{bai2022constitutional}. In these cases, the only human input is a high-level text guideline that informs the model's feedback generation. Collectively, these strategies bridge research in reinforcement learning-based language model training for diverse objectives~\citep{Ranzato2015SequenceLT,paulus2018a,wu2018learning} and general approaches to preference learning from humans \citep{christiano2017deep,kupcsik2018learning}. Despite the attractiveness of utilizing comparative human judgments, fine-tuning LLMs with reinforcement learning poses significant engineering and theoretical hurdles; the current work introduces a principled technique for optimizing relative preferences that eliminates the need for reinforcement learning.

Research on learning policies from preferences outside the linguistic domain has encompassed both bandit and reinforcement learning frameworks, with a variety of methodologies introduced. In contextual bandits, learning is often guided by preferences or action rankings instead of explicit reward signals, leading to the formulation of the contextual dueling bandit (CDB) paradigm \citep{yue2012karmed,dudik2015contextual}. Without direct access to reward values, analyses of CDBs employ the concept of a \textit{von Neumann winner}: a policy that achieves an expected victory rate of at least 50\% when compared to any other policy \citep{dudik2015contextual}. Notably, CDB scenarios generally assume that preference feedback is acquired in an online manner, whereas learning from human-provided preferences frequently relies on a fixed, offline dataset comprising action pairs annotated with preferences \citep{yan2022human}. In a related vein, \textit{preference-based RL} (PbRL) utilizes binary preferences informed by an unknown ‘scoring’ function, as opposed to direct rewards \citep{BusaFekete2014,ruiz2023dueling}. Although a spectrum of PbRL algorithms has been introduced—including those capable of exploiting off-policy preference information—they commonly follow a two-step process: first estimating the underlying scoring or reward function, and subsequently performing optimization based on that estimate \citep{jain2013learning,BusaFekete2014,christiano2017deep,sadigh2017active,kupcsik2018learning}. By contrast, our work advocates for a one-stage approach where policy optimization is conducted directly to align with observed preferences.

\section{Preliminaries}\label{section:prelims}

We summarize the RLHF process outlined by \citeauthor{ziegler2020finetuning}, as well as in subsequent works \citep{stiennon2022learning, bai2022training, ouyang2022training}. This methodology typically unfolds in three main steps: (1) supervised fine-tuning (SFT); (2) collecting preference data and training a reward model; and (3) reinforcement learning-based optimization.

\textbf{SFT}: The RLHF pipeline begins with supervised fine-tuning of a pre-trained language model, utilizing high-quality labeled datasets for target applications such as dialogue or summarization, resulting in the supervised fine-tuned policy $\pi^\text{SFT}$.

\textbf{Reward Modelling Phase}: In the subsequent stage, the SFT policy generates paired outputs $(y_1, y_2)\sim \pi^\text{SFT}(y \mid x)$ in response to prompts $x$. Human annotators then indicate their preference for one of the completions; this is denoted $y_w\succ y_l \mid x$, with $y_w$ and $y_l$ indicating the preferred and non-preferred responses, respectively, among $(y_1, y_2)$. It is posited that such preferences reflect an underlying latent reward function $r^*(y, x)$, which remains unobservable. Several probabilistic models exist to account for human preferences, with the Bradley-Terry (BT) model \cite{bradley1952rankanalysis} being widely used (and more general frameworks like the Plackett-Luce model \citep{plackett1975analysis, luce2012individual} are applicable when multiple rankings are available). According to the BT model, the distribution $p^*$ expressing human preferences is given by:

Given a static dataset of comparison data $\mathcal{D}=\bigl\{x^{(i)}, y_w^{(i)}, y_l^{(i)}\bigr\}_{i=1}^N$ that is drawn from $p^*$, it is standard to model a reward function $r_{\phi}(x, y)$ and to fit the parameters via the principle of maximum likelihood. When cast as a binary classification problem, the optimization objective becomes the negative log-likelihood loss:

\begin{equation}\label{eq:reward_model}
    \mathcal{L}_R(r_{\phi}, \mathcal{D}) = -\mathbb{E}_{(x, y_w, y_l)\sim \mathcal{D}}\bigl[\log \sigma(r_{\phi}(x, y_w)- r_{\phi}(x, y_l))\bigr]
\end{equation}

where $\sigma$ denotes the logistic (sigmoid) function. For language model applications, $r_{\phi}(x, y)$ is typically initialized from the SFT model $\pi^\text{SFT}(y \mid x)$, extending it with an additional linear layer placed atop the last transformer hidden state to yield a scalar reward output \cite{ziegler2020finetuning}. To limit reward variance, prior work has normalized the reward outputs so that $\mathbb{E}_{x,y\sim \mathcal{D}}\left[r_\phi(x, y)\right] = 0$ for any input $x$.

\textbf{RL Fine-Tuning Phase}: In the subsequent RL stage, the learned reward model serves as a feedback mechanism for the language model’s outputs. Building on earlier work~\citep{jaques2017sequence, jaques2020human}, the optimization objective can be written as

\begin{equation}\label{eq:RL}
\max_{\pi_{\theta}}  \mathbb{E}_{x\sim \mathcal{D}, y\sim \pi_{\theta}(y \mid x)}\bigl[r_{\phi}(x, y)\bigr] - \beta\mathbb{D}_{\textrm{KL}}\bigl[\pi_{\theta}(y\mid x)\mid \mid \pi_\text{ref}(y\mid x)\bigr],
\end{equation}

with $\beta$ as a regularization parameter governing the extent to which the policy $\pi_\theta$ can diverge from the reference policy $\pi_\text{ref}$, usually chosen as the original SFT model $\pi^\text{SFT}$. In most implementations, $\pi_\theta$ is initialized to $\pi^\text{SFT}$ as well. The presence of the KL divergence term ensures that the learned policy remains within the support of the data for which the reward model is reliable, and helps to retain output diversity, thus avoiding degeneration into a single dominant response. Due to the inherently discrete outputs of language models, this objective lacks straightforward differentiability and is commonly solved with reinforcement learning techniques. A prevalent method~\citep{ziegler2020finetuning, stiennon2022learning, bai2022training, ouyang2022training} recasts the objective as maximizing ${r(x, y) = r_{\phi}(x, y) -\beta (\log \pi_{\theta}(y\mid x) - \log \pi_\text{ref}(y\mid x))}$ using the PPO algorithm~\cite{schulman2017proximal}.

\section{Direct Preference Optimization}\label{sec:DPO}

Recognizing the difficulties associated with deploying RL algorithms on large-scale problems, such as fine-tuning high-capacity language models, we seek to introduce a more direct approach for policy improvement that operates on preferences themselves. Whereas earlier RLHF pipelines rely on separately training a reward model and then performing RL optimization with respect to this model, our method is grounded in a specific reward model parameterization that permits closed-form computation of the optimal policy, removing the need for an explicit RL loop. Our main contribution is to exploit an analytical relationship mapping rewards to optimal policies, thereby permitting us to translate a loss over rewards into a loss defined directly on policies. This change of variables circumvents the explicit construction and fitting of a separate reward model, but still aligns training with established frameworks for modeling preference data, such as the Bradley-Terry model. Ultimately, this allows the policy network to simultaneously embody both the generative language model and its (implicit) reward structure.

\textbf{DPO Objective Derivation.} Beginning with the standard RL objective expressed in Eq.~\ref{eq:RL} under a general reward function $r$, as established in previous works~\citep{peters2007reinforcement, peng2019advantage, korbak2022reinforcement, go2023aligning}, it is readily demonstrated that the optimal policy under a KL-divergence constraint, as formulated in Eq.~\ref{eq:RL}, adopts the following structure:

\begin{equation}\label{eq:op_policy}
    \pi_r(y\mid x) = \frac{1}{Z(x)}\pi_\text{ref}(y\mid x)\exp\left(\frac{1}{\beta}r(x, y)\right),
\end{equation}

with the normalization term $Z(x) =\sum_{y}\pi_\text{ref}(y\mid x)\exp\left(\frac{1}{\beta}r(x, y)\right)$. A detailed mathematical derivation is provided in Appendix~\ref{app:derivation1}. Despite having an MLE estimator $r_{\phi}$ for the true reward $r^*$, directly computing the partition function $Z(x)$ is computationally intractable~\citep{korbak2022reinforcement, go2023aligning}, limiting the direct application of this form. Nevertheless, it is possible to rearrange Eq.~\ref{eq:op_policy} to recast the reward in terms of its associated optimal policy $\pi_r$, the reference policy $\pi_\text{ref}$, and the partition function $Z(\cdot)$. By applying the logarithm to both sides of Eq.~\ref{eq:op_policy} and performing straightforward manipulations, we arrive at:

\begin{equation}\label{eq:main_eq}
    r(x,y) =\beta \log \frac{\pi_r(y\mid x)}{\pi_\text{ref}(y\mid x)} + \beta \log Z(x).
\end{equation}

Utilizing this transformation for the true reward $r^*$ and its corresponding optimal policy $\pi^*$, we note that the Bradley-Terry preference model considers only the reward gap between two responses: ${p^*(y_1 \succ y_2 \mid x) = \sigma(r^*(x, y_1) - r^*(x, y_2))}$. Substituting Eq.~\ref{eq:main_eq} for $r^*(x, y)$ into the preference formulation Eq.~\ref{eq:bradley-terry}, the partition function terms cancel out, allowing the human preference probability to depend solely on $\pi^*$ and $\pi_\text{ref}$. Accordingly, under the Bradley-Terry framework, the optimal RLHF policy $\pi^*$ fulfills the condition:

\begin{equation}\label{eq:objective}
    p^*(y_1\succ y_2 \mid x)=\frac{1}{1 + \exp\left(\beta \log \frac{\pi^*(y_2\mid x)}{\pi_\text{ref}(y_2\mid x)} - \beta \log \frac{\pi^*(y_1\mid x)}{\pi_\text{ref}(y_1\mid x)}\right)}
\end{equation}

The full derivation appears in Appendix~\ref{app:derivation2}. While Eq.~\ref{eq:objective} employs the Bradley-Terry approach, analogous derivations are possible for broader classes of models, such as Plackett-Luce~\citep{plackett1975analysis, luce2012individual}, as documented in Appendix~\ref{app:plackett_luce_models}.

Given that we now represent human preference probabilities in terms of the optimal policy, rather than through the reward function, we can establish a maximum likelihood training objective for a parameterized policy $\pi_\theta$. In line with reward modeling frameworks (see Eq.~\ref{eq:reward_model}), our policy learning objective is formulated as:

\begin{equation}\label{eq:optimum_model}
    \mathcal{L}_\text{DPO}(\pi_{\theta}; \pi_\text{ref}) = -\mathbb{E}_{(x, y_w, y_l)\sim \mathcal{D}}\left[\log \sigma \left(\beta \log \frac{\pi_{\theta}(y_w\mid x)}{\pi_\text{ref}(y_w\mid x)} - \beta \log \frac{\pi_{\theta}(y_l\mid x)}{\pi_\text{ref}(y

In this formulation, we implicitly learn a reward through an alternative parameterization, such that the resulting optimal policy aligns with $\pi_\theta$. Additionally, our approach is tantamount to fitting a reparameterized Bradley-Terry model, which inherits certain theoretical guarantees—specifically, consistency—provided that the preference data distribution satisfies suitable assumptions \cite{bong2022generalized}. We elaborate further on the theoretical characteristics of DPO and its relationship to prior literature in Section~\ref{sec:theory}.

\textbf{How does the DPO update function?} To gain a mechanistic perspective on DPO, it is instructive to examine the gradient of its objective $\mathcal{L}_\text{DPO}$. The gradient with respect to the parameters $\theta$ is given by:

\begin{multline*}\label{eq:gradient}
    \nabla_\theta \mathcal{L}_\text{DPO}(\pi_\theta;\pi_\text{ref}) = \\ -\beta\mathbb{E}_{(x, y_w, y_l) \sim \mathcal{D}} \bigg[\underbrace{\sigma(\hat{r}_\theta(x, y_l) - \hat{r}_\theta (x, y_w))}_\text{larger impact when reward assignment is erroneous}\bigg[\underbrace{\nabla_\theta\log \pi(y_w \mid x)}_\text{encourage $y_w$} - \underbrace{\nabla_\theta\log\pi(y_l \mid x)}_\text{discourage $y_l$}\bigg]\bigg],
\end{multline*}

where the implicit reward function is defined as $\hat{r}_\theta(x, y) = \beta \log \frac{\pi_\theta(y \mid x)}{\pi_\text{ref}(y \mid x)}$, drawing upon both the language model $\pi_\theta$ and the reference $\pi_\text{ref}$ (further details in Section~\ref{sec:theory}). In essence, the gradient of $\mathcal{L}_\text{DPO}$ incentivizes higher probabilities for preferred outputs $y_w$ and diminishes those for less preferred outputs $y_l$. Crucially, the contribution of each training sample is modulated by the degree to which the implicit reward model $\hat{r}_\theta$ assigns greater value to the non-preferred completion, scaled by $\beta$. This mechanism thus captures both the severity of incorrect reward ranking and the influence of the KL constraint. Empirically, we find this sample weighting to be pivotal, as omitting the weighting factor in a simpler variant often leads to language model collapse (see Appendix Table~\ref{tab:unlikelihood_generations}).

\textbf{DPO overview.} The DPO pipeline can be summarized in the following steps: 1) For each prompt $x$, generate completions $y_1, y_2 \sim \pi_\text{ref}(\cdot \mid x)$, and assign labels according to human preference to build an offline preference dataset $\mathcal{D} = \{x^{(i)}, y_w^{(i)}, y_l^{(i)}\}_{i=1}^N$; 2) Adjust the language model $\pi_\theta$ by minimizing the DPO loss $\mathcal{L}_\text{DPO}$, given $\pi_\text{ref}$, $\mathcal{D}$, and the target $\beta$ parameter. In practice, rather than collecting new human feedback, it is often desirable to utilize publicly available preference datasets. Because these datasets are typically generated using $\pi^\text{SFT}$, we set $\pi_\text{ref} = \pi^\text{SFT}$ when this model is accessible. If $\pi^\text{SFT}$ is not provided, we instead set $\pi_\text{ref}$ by maximizing the likelihood of preferred completions, i.e., we choose ${\pi_\text{ref} = \argmax_{\pi}\mathbb{E}_{x, y_w \sim \mathcal{D}}\left[\log \pi(y_w \mid x)\right]}$. This strategy helps reduce the gap between the unavailable true reference distribution and the $\pi_\text{ref}$ applied in DPO. Details on implementation and hyperparameter selection are elaborated in Appendix~\ref{app:implementation}.

\section{Theoretical Analysis of DPO} This section provides a deeper examination of the DPO framework, theoretical justification for its approach, and discusses the strengths of DPO relative to limitations found in actor-critic methods such as those used in RLHF (for instance, PPO~\cite{schulman2017proximal}).

\label{sec:theory}

\subsection{Your Language Model Is Secretly a Reward Model} DPO eliminates the need for both explicit reward modeling and reinforcement learning by leveraging a single maximum likelihood loss. In fact, the optimization target in Eq. \ref{eq:main_eq} aligns with the Bradley-Terry model when using the reward function $r^*(x, y) = \beta \log\frac{\pi^*_\theta(y \mid x)}{\pi_\text{ref}(y \mid x)}$. The parameters of our model $\pi_{\theta}$ are updated analogously to the reward model approach in Eq. \ref{eq:reward_model}, up to a reparametrization. In what follows, we establish the theoretical foundation for this reparametrization, demonstrate that it does not limit the expressive power of learned reward models, and prove that it supports perfect recovery of the optimal policy. We begin by introducing an equivalence relation for reward functions.

\begin{definition}
Two reward functions $r(x, y)$ and $r'(x, y)$ are considered equivalent if and only if there exists a function $f$ such that ${r(x, y)-r'(x, y) = f(x)}$.
\end{definition}

It is straightforward to verify that this definition indeed describes an equivalence relation, which divides the space of reward functions into equivalence classes. This leads to the following two lemmas:

\begin{lemma}\label{lemma:same_prefrence}
    Within the Plackett-Luce model, including its special case Bradley-Terry, any pair of reward functions belonging to the same equivalence class will produce identical preference distributions.
\end{lemma}

\begin{lemma}\label{lemma:same_policy}
    Any two reward functions from the same equivalence class lead to the same optimal policy when solving the corresponding constrained reinforcement learning problem.
\end{lemma}

The proofs for these statements are uncomplicated and are provided in Appendix \ref{app:lemma1}. The first lemma highlights a well-established issue of under-specification present in the Plackett-Luce model family \cite{plackett1975analysis}. Owing to this ambiguity, it is standard practice to impose supplementary identifiability conditions to secure properties such as unique maximum likelihood estimates in Eq.~\ref{eq:reward_model} \cite{bong2022generalized}. The second lemma indicates that all members of an equivalence class induce the same optimal policy. Thus, for the purposes of our final objective, it suffices to recover any representative reward function from the optimal equivalence class. We establish the following theorem in Appendix~\ref{app:thm1}:

\begin{theorem}\label{thm:main}
    Provided certain regularity conditions, every equivalence class of reward functions compatible with the Plackett-Luce (and, by extension, Bradley-Terry) models admits a representation via the reparameterization ${r(x, y) = \beta \log \frac{\pi(y\mid x)}{\pi_\text{ref}(y\mid x)}}$, for some model $\pi(y\mid x)$ and given reference model $\pi_\text{ref}(y \mid x)$.
\end{theorem}

\begin{sproof}
    Take any reward function $r(x, y)$, and let the associated optimal policy be $\pi_r(y \mid x)$, as defined in Eq.~\ref{eq:op_policy}. We claim that a member of the equivalence class containing $r$ can be constructed in the desired reparameterized form. Let us define a projection operator $f$ as
\begin{equation}
    f(r; \pi_\text{ref}, \beta)(x, y) = r(x, y) - \beta\log\sum_{y}\pi_\text{ref}(y\mid x)\exp\left(\frac{1}{\beta}r(x, y)\right)
\end{equation}
This operation amounts to normalizing the reward function by subtracting the logarithm of the partition function corresponding to $\pi_r$, which only depends on $x$. Therefore, $f(r; \pi_\text{ref}, \beta)(x, y)$ lies within the equivalence class defined by $r(x, y)$. Substituting the right-hand side of Eq.~\ref{eq:main_eq} (valid for all reward functions) into $r$, we have $f(r; \pi_\text{ref}, \beta)(x, y) = \beta \log \frac{\pi_r(y\mid x)}{\pi_\text{ref}(y\mid x)}$. This shows that the projection $f$ yields a reward function in the form stated, ensuring the generality of the proposed reparameterization within the equivalence class.
\end{sproof}

An alternative interpretation of Theorem~\ref{thm:main} is that it uniquely determines the specific reward function chosen by the DPO reparameterization within each equivalence class—namely, the one satisfying

\begin{equation}\label{eq:lag_p}
     \sum_{y}\underbrace{\pi_\text{ref}(y\mid x)\exp\left(\frac{1}{\beta}r(x, y)\right)}_{=\pi(y\mid x)\text{, using Thm.~\ref{thm:main} reparam.}} = 1,
\end{equation}

This ensures that $\pi(y\mid x)$ defines a valid probability distribution, with non-negative values summing to one. Furthermore, based on Eq.~\ref{eq:op_policy}, we observe that Eq.~\ref{eq:lag_p} corresponds to the partition function associated with the optimal policy generated by the reward function $r(x, y)$. The core insight underlying the DPO method is its ability to introduce constraints into the under-determined Plackett-Luce (and, specifically, Bradley-Terry) family of preference models. By doing so, we maintain the breadth of representable reward models, while also ensuring that the optimal policy described in Eq.~\ref{eq:op_policy} remains analytically manageable for every prompt $x$.

\subsection{Instability of Actor-Critic Algorithms} 

Our analytical framework also helps elucidate sources of instability in typical actor-critic approaches adopted for RLHF, such as PPO. We follow the RLHF methodology and concentrate on the RL fine-tuning component as discussed in Section \ref{section:prelims}. There is a notable connection with the control as inference paradigm \cite{levine2018reinforcement} when dealing with the constrained RL formulation introduced in \ref{eq:RL}. Suppose we have a model $\pi_{\theta}(y\mid x)$ with parameters $\theta$, and we aim to minimize $\mathbb{D}_{\text{KL}}[\pi_{\theta}(y|x) \mid \mid \pi^*(y\mid x)]$, where $\pi^*$ is the optimal policy derived from Eq.~\ref{eq:optimum_model} and is shaped by the reward function $r_{\phi}(y, x)$. Manipulating these terms yields the following optimization objective:

\begin{equation}\label{eq:AC}
    \max_{\pi_{\theta}}\mathbb{E}_{\pi_{\theta}(y\mid x)}\bigg[\underbrace{r_{\phi}(x, y) -\beta\log\sum_{y}\pi_\text{ref}(y\mid x)\exp\left(\frac{1}{\beta}r_{\phi}(x, y)\right)}_{f(r_{\phi}, \pi_\text{ref}, \beta)} - \underbrace{\beta\log\frac{\pi_{\theta}(y\mid x)}{\pi_\text{ref}(y\mid x)}}_{\text{KL}}\bigg]
\end{equation}

The objective function under consideration aligns with those optimized in earlier studies 
\citep{ziegler2020finetuning, stiennon2022learning, bai2022training, ouyang2022training}, which employ the DPO-equivalent reward framework for the reward parameterization $r_{\phi}$. Within this context, the normalization factor in $f(r_{\phi}, \pi_\text{ref}, \beta)$ can be seen as the soft value function associated with the reference policy $\pi_\text{ref}$. Although this normalization term does not impact the optimal solution, omitting it may result in a policy gradient with elevated variance, potentially destabilizing the learning dynamics. One possible strategy to address the normalization is to introduce a learned value function; however, this too presents optimization challenges. Traditionally, previous approaches have dealt with normalization by referencing a human completion baseline, effectively using a single-sample Monte-Carlo approximation for the normalizing factor. In contrast, the DPO reparameterization introduces a reward function formulation that obviates the need for any baseline normalization.

\section{Experiments}
Here, we conduct empirical assessments of DPO with regard to training policies directly from preference data. Our investigation begins with a controlled text-generation scenario, where we examine how effectively DPO balances reward maximization and KL-divergence minimization with respect to the reference policy, as compared to standard preference learning approaches such as PPO. Subsequently, we test DPO on more complex models and RLHF benchmarks, including summarization and dialogue tasks. Remarkably, DPO generally achieves performance on par with or superior to robust baselines—including RLHF with PPO and best-of-$N$ sampling strategies under learned reward models—even with minimal hyperparameter tuning. Prior to presenting our empirical findings, we detail our experimental methodology; comprehensive information can be found in Appendix~\ref{app:exp_details}.

\textbf{Tasks.} In this work, we assess performance across three open-domain text generation tasks. Across all tasks, models are trained to learn a policy from a preference dataset denoted as $\mathcal{D}=\bigl\{x^{(i)}, y_w^{(i)}, y_l^{(i)}\bigr\}_{i=1}^N$. For \textbf{controlled sentiment generation}, $x$ represents a prompt extracted from a movie review in the IMDb dataset \cite{maas-EtAl:2011:ACL-HLT2011}, and the policy's objective is to generate $y$ exhibiting positive sentiment. To systematically evaluate performance in this domain, we automatically construct preference pairs for each generated sample by leveraging a pre-trained sentiment classifier; these are constructed such that $p(\text{positive}\mid x,y_w)>p(\text{positive}\mid x,y_l)$. In the SFT baseline, we perform fine-tuning of GPT-2-large on the training portion of IMDb reviews until satisfactory convergence is achieved (see App~\ref{app:sentiment_details} for additional methodological specifics). For the \textbf{summarization} task, $x$ consists of a Reddit forum post, while the policy must produce a concise summary $y$ that distills the post's core arguments. In line with prior literature, we make use of the Reddit TL;DR dataset \citep{volske-etal-2017-tl}, incorporating human preference annotations curated by \citeauthor{stiennon2022learning}. The SFT baseline is obtained by fine-tuning on summaries authored by humans for Reddit forum content,\footnote{\url{https://huggingface.co/CarperAI/openai_summarize_tldr_sft}} implemented with the TRLX \citep{leandro_von_werra_2023_7790115} package for RLHF. The dataset of human preferences was previously collected by \citeauthor{stiennon2022learning} from a set of generations produced by a distinct but equivalently fine-tuned SFT model. Lastly, in \textbf{single-turn dialogue}, $x$ is a user message that can range from scientific queries to personal advice requests. The model's policy is to respond to the query with an answer $y$ that is both helpful and engaging. This task employs the Anthropic Helpful and Harmless dialogue dataset \citep{bai2022training}, consisting of 170,000 human–AI dialogue exchanges. Each dialogue ends with two assistant responses generated by a large, unspecified language model, along with a human-annotated label indicating the favored response. Unlike other settings, there is no pre-trained SFT model available; thus, we construct the SFT baseline by further training a standard language model solely on the preferred completions.

\textbf{Evaluation.} Our empirical assessment employs two distinct methodologies. To specifically investigate how well each algorithm maximizes the reward objective under constraints, we examine their performance in a controlled sentiment generation scenario by mapping out the Pareto frontier between realized reward and KL-divergence relative to a reference policy; this evaluation is feasible because we possess the true reward signal through an accessible sentiment classifier. Contrastingly, in practical applications where the actual reward function is unavailable, we resort to measuring algorithms' \textit{win rate} against a baseline using GPT-4 as an automated stand-in for human raters. This proxy is applied to judge the quality of generated summaries and the helpfulness of single-turn dialogue responses. For summarization, our baseline is the human-generated reference summaries from the test split; for dialogue, it is the response marked as preferred within the test set. Although prior research suggests that LMs might surpass standard automated metrics in evaluation fidelity \citep{Chen2023ExploringTU}, we perform a human study (see Sec.~\ref{sec:human-judgments}) to further support our selection of GPT-4 for this role. Our findings indicate a high degree of concordance between GPT-4 and human annotations, frequently matching or exceeding the level of agreement found between different human evaluators.

\textbf{Methods.} Alongside DPO, we compare a suite of established strategies for aligning language model outputs with human preferences. As straightforward baselines, we test zero-shot prompting using \textbf{GPT-J} \citep{gpt-j} for summarization and 2-shot prompting with \textbf{Pythia-2.8B} \citep{biderman2023pythia} in dialogue generation. Additionally, our benchmarks include the \textbf{SFT} model and \textbf{Preferred-FT}, where the model is supervised via fine-tuning on the selected output $y_w$, originating either from the SFT model (in sentiment-controlled and summarization contexts) or from a generic LM (for dialogue tasks). Another form of pseudo-supervised training, \textbf{Unlikelihood}~\citep{welleck2019neural}, seeks to enhance the policy by maximizing the likelihood assigned to $y_w$ while penalizing the likelihood of $y_l$, modulated by a tunable unlikelihood coefficient $\alpha\in[0,1]$. We further include \textbf{PPO} \citep{schulman2017proximal} leveraging a reward model learned from preferences, and the oracle-like \textbf{PPO-GT}, which uses the actual reward function available in the controlled sentiment domain. In our sentiment experiments, PPO-GT is represented in two forms: a standard implementation \cite{leandro_von_werra_2023_7790115} and an enhanced version incorporating reward normalization and additional hyperparameter tuning to boost effectiveness (these modifications are similarly applied to PPO with learned rewards). We also evaluate the \textbf{Best of $N$} strategy: by sampling $N$ candidate completions from the SFT model (or Preferred-FT for dialogue) and selecting the one with the highest predicted reward, as determined by a preference-trained model. While this approach generally delivers strong results by isolating the reward model’s performance from PPO training dynamics, its computational expense—necessitating $N$ samples per test query—renders it impractical at even moderate values of $N$.

\subsection{How well can DPO optimize the RLHF objective?}

The reward maximization objective in standard RLHF approaches is constrained by KL-divergence, ensuring that the policy achieves high reward without straying excessively from the reference policy. Consequently, evaluations of different algorithms should consider both attained reward and the corresponding KL divergence; an algorithm that obtains marginally better rewards at the cost of substantially increased KL may not be advantageous. Figure~\ref{fig:frontier-tldr-main} illustrates the trade-off between reward and KL across various methods within the sentiment analysis scenario. For each algorithm, we conduct several training sessions, each time adjusting a distinct policy conservativeness hyperparameter (specifically: target KL $\in\{3,6,9,12\}$ for PPO, $\beta \in \{0.05,0.1,1,5\}$ for DPO, $\alpha\in\{0.05,0.1,0.5,1\}$ for unlikelihood, and random seeds for preferred-FT), accumulating 22 total runs. At intervals of every 100 training steps, up to convergence, we assess every resulting policy on a designated set of test prompts, measuring both the mean reward according to the ground-truth reward function and the average sequence-level KL\footnote{This refers to the aggregate of the KL-divergence computed at each time step.} with the reference policy, denoted by $\text{KL}\left(\pi\mid \mid \pi_\text{ref}\right)$. The findings indicate that DPO delineates a superior frontier, consistently obtaining higher rewards at lower KL values compared to other approaches. This observation is striking for several reasons: firstly, while both DPO and PPO target the same objective, DPO achieves superior efficiency, outpacing PPO across all reward/KL trade-offs; secondly, DPO surpasses PPO even under circumstances where PPO is granted access to true reward signals (PPO-GT).

\subsection{Can DPO scale to real preference datasets?}
\label{sec:dpo-real-datasets}
We proceed to assess DPO's fine-tuning efficacy on tasks in summarization and single-turn dialogue. In the domain of summarization, established automatic metrics like ROUGE often show limited alignment with human judgments~\citep{stiennon2022learning}, and previous research has demonstrated that leveraging PPO-based fine-tuning on human preference data results in improved summary quality. Our evaluation involves generating model outputs on the TL;DR summarization dataset's test partition and calculating the mean win rate versus reference outputs in this set. Sampling is performed for each approach across temperatures from 0.0 to 1.0, with the resulting win rates depicted in Figure~\ref{fig:frontier-tldr-main} (right). DPO, PPO, and Preferred-FT are all initialized from the same GPT-J SFT checkpoint\footnote{\url{https://huggingface.co/CarperAI/openai_summarize_tldr_sft}}. Our analysis shows that DPO attains a win rate near 61\% at a sampling temperature of 0.0, outperforming PPO, which reaches about 57\% under optimal conditions. Additionally, DPO secures a superior maximal win rate relative to the strongest $N$ baseline. It is important to highlight that the $\beta$ hyperparameter for DPO was not subject to extensive tuning; therefore, reported outcomes may not reflect the method’s best possible performance. Notably, DPO exhibits greater robustness to sampling temperature changes than PPO, whose results may revert to the base GPT-J level as temperature increases. Meanwhile, Preferred-FT does not yield appreciable gains over the initial SFT model. Finally, a direct human evaluation of DPO and PPO, described in Section~\ref{sec:human-judgments}, finds that DPO outputs sampled at temperature 0.25 were chosen 58\% of the time over PPO outputs sampled at the same temperature.

For single-turn dialogue, we assess the various approaches using the subset of the Anthropic HH test set \citep{bai2022training} that involves a single exchange between human and assistant. To calculate the win rates for the different techniques, GPT-4 is used to compare model outputs against the preferred completions in the test set, which serve as the ground truth. Since a standard SFT model is not available for this benchmark, our process begins with a pre-trained Pythia-2.8B. We apply Preferred-FT to train a reference model on the preferred completions, ensuring that output samples remain within the intended distribution, followed by DPO training. Additionally, we benchmark against the strongest performer from 128 Preferred-FT completions (as our analysis—see Appendix Figure~\ref{fig:best-of-n}—indicates the Best of $N$ metric saturates at $N=128$), as well as a 2-shot prompt configuration of the unmodified Pythia-2.8B. Across these experiments, DPO is found to achieve equal or superior results for each method's optimal temperature. We also assess an RLHF model fine-tuned with PPO on the Anthropic HH data \footnote{\url{https://huggingface.co/reciprocate/ppo_hh_pythia-6B}} using a recognized implementation \footnote{\url{https://github.com/CarperAI/trlx/tree/main/examples/hh}}, but are unable to identify a prompting strategy or sampling temperature that surpasses the performance of the original Pythia-2.8B. In view of our findings on TL;DR and given that both approaches maximize the same reward objective, we regard the Best of 128 strategy as an approximate stand-in for PPO-level results. Overall, DPO emerges as the only method that offers computational efficiency while still outperforming the preferred completions on the Anthropic HH set, matching or exceeding the heavily resource-dependent Best of 128. Finally, as illustrated in Figure~\ref{fig:dialogue-main}, DPO attains its optimal performance in relatively few training steps.

\subsection{Generalization to a new input distribution}

To further analyze the robustness of PPO and DPO under distributional shifts, we assess both policies—trained on the Reddit TL;DR summarization dataset—by applying them to an out-of-distribution dataset: news articles from the CNN/DailyMail test set \citep{nallapati-etal-2016-abstractive}. We employ the best sampling temperatures identified on the TL;DR data (0 and 0.25, respectively). The findings, shown in Table~\ref{tab:ood}, indicate that DPO maintains a notable performance advantage over PPO on this unfamiliar distribution. GPT-4 win rates were determined relative to the reference summaries in CNN/DailyMail, following the same GPT-4 (C) prompt used for Reddit TL;DR, with a simple substitution of “forum post” for “news article.” These results provide preliminary support that DPO’s learned policies exhibit generalization capabilities on par with those of PPO, even though DPO does not leverage the extra unlabeled Reddit TL;DR prompts available to PPO.

\subsection{Validating GPT-4 judgments with human judgments}
\label{sec:human-judgments}
We implement a human evaluation to assess the reliability of GPT-4’s scoring, drawing from the Reddit TL;DR summarization experiment and two distinct GPT-4 prompting strategies. The \textbf{GPT-4 (S)} (simple) prompt requests a selection of the summary that more effectively encapsulates the core information of a post. In contrast, the \textbf{GPT-4 (C)} (concise) prompt also considers summary conciseness; this is motivated by our observation that, under the simple prompt, GPT-4 is inclined to favor summaries that are longer and more repetitive than those preferred by humans. Complete versions of the prompts can be found in Appendix~\ref{app:prompts}. To capture a range of summary qualities, we conduct three evaluations, comparing the best-performing (DPO, temp. 0.25), worst-performing (PPO, temp. 1.0), and a mid-level (SFT, temp. 0.25) system to the best PPO model (greedy sampling). Analysis reveals that for both prompts, GPT-4’s preferences are generally as consistent with human judgments as human annotators are with each other; this implies GPT-4 can serve as a practical stand-in for direct human comparison (though multiple human annotations are only available for the DPO and PPO-1 assessments due to limited resources). The \textbf{GPT-4 (C)} prompt tends to yield win rates most aligned with human outcomes, prompting us to adopt this setup for the central evaluations in Section~\ref{sec:dpo-real-datasets}. Further details of the human evaluation protocol, including the interface used and participant roster, are included in Appendix~\ref{app:human-study}.

\section{Discussion}
Preference-based learning presents a robust and scalable approach for developing effective and aligned language models. In this work, we have proposed DPO, a straightforward methodology for leveraging preferences to train language models without relying on reinforcement learning techniques. Instead of re-framing preference learning to fit conventional RL architectures, DPO formulates a direct correspondence between language model policies and reward functions. This framework allows for the training of language models to align with human preferences using a basic cross-entropy objective, all without incorporating reinforcement learning or compromising generality. Notably, DPO achieves comparable or superior results to established RLHF strategies—such as those based on PPO—even when hyperparameter tuning is minimal, significantly lowering the entry threshold for developing language models attuned to human preferences.

\textbf{Limitations \& Future Work.} Several avenues merit deeper investigation based on our findings. One critical open question concerns the out-of-distribution generalization capabilities of DPO policies relative to those trained with explicit reward functions; preliminary evidence points toward DPO achieving generalization on par with PPO-derived models, yet further, more systematic evaluation is required. Moreover, investigating whether self-labeling approaches using DPO policies can efficiently exploit unlabeled prompts is an area for future exploration. Another important aspect is understanding how reward over-optimization emerges within the direct preference optimization context, and whether the slight performance dip noted in Figure~\ref{fig:dialogue-main}-right serves as a case in point. Although our experiments cover models with up to 6B parameters, future research could extend to much larger, state-of-the-art architectures. On the assessment front, we observe that GPT-4’s win rate metrics are susceptible to prompt wording; identifying optimal strategies for obtaining accurate evaluations from automated judges remains an open challenge. Lastly, DPO’s utility is not limited to language models—it offers potential for training generative models across various modalities guided by preference data.